\documentclass[9pt]{extarticle}
\usepackage[a4paper,margin=13mm]{geometry}
\usepackage[T1]{fontenc}
\usepackage{lmodern}
\usepackage{booktabs}
\usepackage{longtable}
\usepackage{array}
\usepackage{xcolor}
\usepackage{colortbl}
\usepackage{parskip}
\usepackage{titlesec}

\definecolor{headbg}{RGB}{38,66,105}
\definecolor{ctrl}{RGB}{176,96,42}
\definecolor{cmd}{RGB}{122,79,176}
\definecolor{data}{RGB}{42,143,134}
\renewcommand{\arraystretch}{1.15}
\titleformat{\section}{\large\bfseries\color{headbg}}{\thesection}{0.5em}{}
\titlespacing*{\section}{0pt}{12pt}{4pt}
\titleformat{\subsection}{\normalsize\bfseries}{}{0em}{}
\titlespacing*{\subsection}{0pt}{7pt}{2pt}

\newcolumntype{N}{>{\ttfamily\footnotesize\raggedright\arraybackslash}p{3.9cm}}
\newcolumntype{D}{>{\footnotesize\centering\arraybackslash}p{0.9cm}}
\newcolumntype{W}{>{\ttfamily\footnotesize\centering\arraybackslash}p{2.5cm}}
\newcolumntype{R}{>{\footnotesize\centering\arraybackslash}p{1.5cm}}
\newcolumntype{X}{>{\footnotesize\raggedright\arraybackslash}p{7.1cm}}

\newenvironment{ports}{%
  \footnotesize\begin{longtable}{N D W R X}\toprule
  \textbf{Signal} & \textbf{Dir} & \textbf{Width} & \textbf{Ref clk} & \textbf{Meaning}\\ \midrule\endhead}
  {\bottomrule\end{longtable}}

\begin{document}
\begin{center}
{\LARGE\textbf{DFI 5.2 Port \& Drive-Logic Reference}}\\[2pt]
{\small Post-phaser DFI adapter, one channel, \textbf{DDR5} (DDR1--4 later). Source: DFI 5.2
(Cadence, Oct 2024). Config: \texttt{N\_RANKS}=2, DQ=32b sub-channel, \texttt{BL}=16,
\texttt{GEAR}$\in\{1,2,4\}$=phases/DFI-clk, CA=14b. Dir = from MC (out) or from PHY (in).}
\end{center}

\section{Clocking --- one DFI clock, three reference domains}
\textbf{Every DFI signal is clocked by the single \texttt{dfi\_clk}} (DFI 5.2 Table 2). There is no
separate ``DDR clock'' the MC logic sits in --- the PHY generates the DRAM CK behind the phase
interface. Each signal is only \emph{associated} with a reference domain for its timing/phase
relationship:
\begin{itemize}\itemsep1pt
\item \textcolor{cmd}{\textbf{Command clock}}: \texttt{dfi\_address, dfi\_cs, dfi\_cke, dfi\_odt,
  dfi\_reset\_n, dfi\_parity\_in, dfi\_alert\_n, dfi\_2n\_mode, dfi\_cs\_geardown,
  dfi\_dram\_clk\_disable} (+ legacy \texttt{dfi\_act\_n/ras\_n/cas\_n/we\_n/bank/bg/cid}).
\item \textcolor{data}{\textbf{Data clock}}: \texttt{dfi\_wrdata*, dfi\_rddata*, dfi\_wck\_*}.
\item \textcolor{ctrl}{\textbf{Control clock}}: \texttt{dfi\_init\_*, dfi\_*\_freq\_ratio,
  dfi\_freq\_fsp, dfi\_frequency, dfi\_ctrlupd\_*, dfi\_phyupd\_*, dfi\_error*, dfi\_lp\_*}.
\end{itemize}
\textbf{Frequency ratio (GEAR):} in a 1:N system each signal is replicated per phase with a
\texttt{\_pN} suffix (\texttt{\_p0} may drop the suffix); N=GEAR command/data phases ride one
\texttt{dfi\_clk} edge. \textbf{Baseline (STAGE 4): matched, \texttt{dfi\_clk == mc\_clk}} --- the
whole post-phaser adapter is ONE clock domain (\texttt{mc\_clk}); the reference domains only tell
you which phase a bus beat lands on. So your instinct is right that the enables are ``DDR-clock''
by reference, but physically the MC drives them on \texttt{dfi\_clk}, per phase.

\textbf{The single clock is the SLOW clock (\texttt{dfi\_clk = CK/GEAR}); the N$\times$ is SPATIAL,
not a faster clock.} At 1:4 \texttt{dfi\_clk} runs at CK/4 and \emph{each edge drives 4 CA phases in
parallel} (\texttt{dfi\_address\_p0..p3}, \texttt{dfi\_cs\_p0..p3}); the PHY serializes those 4 onto
the DRAM CA at full CK. So the MC \emph{never} runs at CK --- it runs at CK/GEAR and emits a whole
bundle's worth of CA/enables per cycle:
\begin{itemize}\itemsep0pt
\item 1:1 $\rightarrow$ 1 phase/edge (\texttt{dfi\_clk}=CK)
\item 1:2 $\rightarrow$ 2 phases/edge (\texttt{dfi\_clk}=CK/2)
\item 1:4 $\rightarrow$ 4 phases/edge (\texttt{dfi\_clk}=CK/4)
\end{itemize}
This is exactly the phaser's N-winner bundle (STAGE 4). A DDR5 command is 2-UI, so it occupies 2 of
the N phase slots; at 1:4 that means up to 2 two-cycle commands per edge (\texttt{[C C A A]}) or a
mix with 1-cycle PRE (fig\_55). \texttt{DFI\_CMD\_ENCODE} builds the N phase-slots; the PHY does the
CK-rate serialization.

\section{\textcolor{cmd}{Command interface} (Command-clk ref, all \texttt{\_pN} per phase)}
\begin{ports}
dfi\_address(\_pN) & out & 14 & cmd & CA bus. DDR5 command is CA-encoded here (2-UI cmds span 2 beats); MC preserves bit order.\\
dfi\_cs(\_pN)      & out & N\_RANKS=2 & cmd & Chip select (active low), one per rank/CS. Asserted on the command's phase(s).\\
dfi\_cke(\_pN)     & out & N\_RANKS=2 & cmd & Clock enable per rank (power-down / self-refresh entry-exit).\\
dfi\_odt(\_pN)     & out & N\_RANKS=2 & cmd & On-die-termination control per rank.\\
dfi\_reset\_n(\_pN)& out & 1 & cmd & DRAM reset (init).\\
dfi\_2n\_mode(\_pN)& out & 1 & cmd & 2N command timing (hold CA $\ge$1 cmd-clk after CS for DDR5 w/o RCD).\\
dfi\_cs\_geardown(\_pN)& out & 1 & cmd & DDR5 CS geardown mode select.\\
dfi\_dram\_clk\_disable(\_pN)& out & per-clk & cmd & Gate the DRAM clock (clock-stop / low power).\\
dfi\_parity\_in(\_pN)& out & 1 & cmd & CA parity forwarded to PHY (tparin\_lat after cmd).\\
dfi\_alert\_n(\_pN)& \textbf{in} & 1 & cmd & PHY $\rightarrow$ MC alert (CA-parity / write-CRC error), tphy\_paritylat.\\
dfi\_act\_n/ras\_n/cas\_n/we\_n(\_pN)& out & 1 each & cmd & \emph{Legacy DDR3/4} discrete command encode. Unused for DDR5 (CA on dfi\_address).\\
dfi\_bank/bg/cid(\_pN)& out & 2/3/? & cmd & \emph{Legacy} DDR3/4 bank/BG; \texttt{cid} = chip-ID for 3DS.\\
\end{ports}

\section{\textcolor{data}{Write data interface} (Data-clk ref, \texttt{\_pN} per phase)}
\begin{ports}
dfi\_wrdata(\_pN)   & out & 2$\cdot$DQ=64 & data & Write data, driven tphy\_wrdata cycles AFTER dfi\_wrdata\_en, for as long as en held.\\
dfi\_wrdata\_en(\_pN)& out & per-slice & data & Marks valid wrdata/mask/crc window. \textbf{Generated by the WL launcher} (tphy\_wrlat from CAS-W).\\
dfi\_wrdata\_mask(\_pN)& out & 2$\cdot$DQ/8=8 & data & Byte mask (DM) or write-DBI. Bit$k$ masks byte$k$ of dfi\_wrdata.\\
dfi\_wrdata\_cs(\_pN)& out & PHY\_RANK\_W & data & Which rank the write data targets (host-timing compensation).\\
dfi\_wrdata\_crc(\_pN)& out & DATA\_W/8 & data & Link CRC per byte-lane, same timing as dfi\_wrdata.\\
dfi\_wrdata\_ecc(\_pN)& out & (opt) & data & Write-data link ECC (if enabled).\\
\end{ports}

\section{\textcolor{data}{Read data interface} (Data-clk ref)}
\begin{ports}
dfi\_rddata(\_wN)    & \textbf{in} & 2$\cdot$DQ=64 & data & Read data from PHY (\texttt{\_wN} = per read-word phase).\\
dfi\_rddata\_en(\_pN)& out & per-slice & data & MC tells PHY the read-capture window. \textbf{Generated by the RL line} (tphy\_rdlat from RD).\\
dfi\_rddata\_valid(\_wN)& \textbf{in} & per-slice & data & PHY $\rightarrow$ MC: this beat of dfi\_rddata is valid (drives RD-capture into RD\_SRAM).\\
dfi\_rddata\_cs(\_pN)& out & PHY\_RANK\_W & data & Rank targeted for the read (host-timing).\\
dfi\_rddata\_dnv(\_wN)& \textbf{in} & per-slice & data & Data-not-valid marker within a returned word.\\
dfi\_rddata\_dbi(\_wN)& \textbf{in} & DATA\_W/8 & data & Read DBI (if enabled).\\
dfi\_rddata\_crc(\_wN)& \textbf{in} & DATA\_W/8 & data & Read link CRC.\\
\end{ports}

\section{\textcolor{ctrl}{Status / Update / Error / Low-power} (Control-clk ref, single-phase)}
\begin{ports}
dfi\_init\_start     & out & 1 & ctrl & MC starts init / requests a frequency change.\\
dfi\_init\_complete  & \textbf{in} & 1 & ctrl & PHY: init done / freq-change ack. Gates the DFI-mux (init$\rightarrow$scheduler).\\
dfi\_cmd\_freq\_ratio& out & 2 & ctrl & Command frequency ratio (1:1/1:2/1:4) = GEAR select.\\
dfi\_data\_freq\_ratio& out & 2 & ctrl & Data frequency ratio.\\
dfi\_freq\_fsp       & out & 1--2 & ctrl & Frequency set-point (which MR timing set). Changes only at init / freq-change.\\
dfi\_frequency       & out & W & ctrl & Target operating frequency code.\\
dfi\_ctrlupd\_req    & out & 1 & ctrl & MC requests a PHY update window.\\
dfi\_ctrlupd\_ack    & \textbf{in} & 1 & ctrl & PHY grants the ctrl update.\\
dfi\_phyupd\_req     & \textbf{in} & 1 & ctrl & PHY requests an update (e.g. periodic retrain).\\
dfi\_phyupd\_ack     & out & 1 & ctrl & MC grants the PHY update.\\
dfi\_phyupd\_type    & \textbf{in} & 2 & ctrl & Kind of PHY update.\\
dfi\_error          & \textbf{in} & 1 & ctrl & PHY error strobe.\\
dfi\_error\_info    & \textbf{in} & W$\times$4 & ctrl & Error source (valid only while dfi\_error).\\
dfi\_lp\_ctrl\_req/ack/wakeup & out/in/out & 1 & ctrl & Low-power control-path handshake.\\
dfi\_lp\_data\_req/ack/wakeup & out/in/out & 1 & ctrl & Low-power data-path handshake.\\
\end{ports}

\section{Post-phaser drive logic --- \textbf{encode after the phaser} (decision)}
\textbf{Question:} generate the DFI enables by (a) a DFI-encode stage after the phaser, or (b) let
the scheduler drive every enable? \textbf{Decision: (a) encode after the phaser.} The scheduler /
phaser output is legality + placement (\emph{what} \& \emph{which phase}); DFI framing (\emph{the
exact per-phase enables and CA beats}) is a separate encode. Reasons:
\begin{itemize}\itemsep1pt
\item The scheduler must stay selection-only; loading it with per-phase \texttt{dfi\_cs}/CA-beat and
  \texttt{wrdata\_en}/\texttt{rddata\_en} timing would couple two unrelated concerns and bloat it.
\item The enables are \textbf{deterministic delays off the CAS commit}, already owned by the data
  path (STAGE 5): \texttt{dfi\_wrdata\_en} from the WL launcher (tphy\_wrlat), \texttt{dfi\_rddata\_en}
  from the RL line (tphy\_rdlat). They are not scheduling decisions.
\item DDR5 CA is 2-UI --- the encoder is the natural place to spread one command over 2 beats and
  stamp \texttt{dfi\_cs} on the right phase(s).
\end{itemize}

\subsection{Blocks after the phaser (all in \texttt{dfi\_clk == mc\_clk})}
\begin{itemize}\itemsep1pt
\item \texttt{DFI\_CMD\_ENCODE} (comb + 1 reg): phaser bundle $\rightarrow$ per-phase
  \{\texttt{dfi\_address\_pN}, \texttt{dfi\_cs\_pN}\}. Maps the DRAM command to the DDR5 CA-bus beats
  and asserts \texttt{dfi\_cs} on the command's phase; fills idle phases with DES (CS high).
\item \texttt{WL\_LAUNCH} (shift line, STAGE 5): at CAS-W commit, arm a tphy\_wrlat delay; on expiry
  drive \texttt{dfi\_wrdata\_en\_pN}, then stream \texttt{dfi\_wrdata\_pN}+\texttt{dfi\_wrdata\_mask\_pN}
  (+CRC) for the burst from \texttt{WD\_SRAM}.
\item \texttt{RD\_CAP} (shift line): at RD commit, arm tphy\_rdlat; assert \texttt{dfi\_rddata\_en\_pN};
  capture \texttt{dfi\_rddata\_wN} on \texttt{dfi\_rddata\_valid\_wN} into \texttt{RD\_SRAM}.
\item \texttt{DFI\_MUX} (init override): selects init-FSM vs scheduler-encoder as CA driver;
  \texttt{sel = dfi\_init\_complete} (the only hard override, STAGE 9).
\item \texttt{ALERT/ERR monitor}: watch \texttt{dfi\_alert\_n} (write-CRC/parity retry window) and
  \texttt{dfi\_error} (stub, STAGE 13/20).
\end{itemize}
\textbf{Sideband (Control-clk):} \texttt{dfi\_init\_*}, \texttt{dfi\_*\_freq\_ratio},
\texttt{dfi\_ctrlupd\_*}, \texttt{dfi\_phyupd\_*} are driven by the maintenance / freq-change FSM,
not the per-command encode.

\section{DDR1--4 portability (to fill later)}
\textbf{Parametric (reload widths / values):} GEAR, DFI Data/Address width, \texttt{N\_RANKS},
CS/CKE/ODT widths, timing params (tphy\_wrlat, tphy\_rdlat, tphy\_wrdata). \textbf{Structural:}
DDR3/4 use the \emph{discrete} command encode (\texttt{dfi\_act\_n/ras\_n/cas\_n/we\_n} +
\texttt{dfi\_bank/bg}) instead of the DDR5 CA-bus-on-\texttt{dfi\_address}, so \texttt{DFI\_CMD\_ENCODE}
gets a per-gen encoder variant; DDR5-only signals (\texttt{dfi\_cs\_geardown}, WCK, MRDIMM mux,
per-slice read valid) drop out. Clock model (single \texttt{dfi\_clk}, 3 ref domains) is identical
across all gens.

\end{document}
